\section{Relationship Between BLEU and Execution-Based Evaluation}
%szb: ehhez szükgésges lenne statisztikai teszteket végezni: mekkora a korreláció? Ez python-ban nagyon egybszerű, csak egy pár sor kód.

The experimental results show that BLEU score and execution-based evaluation
are related, but the relationship between them is not always straightforward.
In general, models with higher average BLEU scores also achieved more accepted
solutions in CodeEditorBench. For example, the fine-tuned Qwen model obtained
the highest BLEU score and also produced the highest number of accepted
translations. Similarly, the fine-tuned StarCoder2 model improved both in BLEU
score and execution-based performance compared to its base version.

However, the results also reveal that higher BLEU scores do not always imply
better functional correctness. In several cases, generated code achieved a
relatively high BLEU score while still failing to compile or producing incorrect
outputs during testing. This happens because BLEU measures token-level overlap
between the generated code and the reference solution, but does not consider
whether the translated code is executable or semantically correct.

This effect can also be observed in the difference between compilation errors
and wrong answers. Fine-tuned models produced fewer compilation errors and more
accepted solutions, but they also produced more wrong answers. This suggests
that fine-tuning often improved syntactic correctness enough for the code to
compile, but some translations still failed to preserve the exact logic of the
original program.

The weak correlation between BLEU and semantic correctness has already been %szb: cite missing
observed in previous research. Studies have shown that BLEU often fails to
reflect the true quality of translated code because multiple implementations can
be functionally equivalent while having very different token sequences.
Similarly, code with high lexical similarity may still contain semantic or
logical errors that prevent correct execution.


The results of this thesis therefore support the idea that execution-based
evaluation is more reliable than BLEU for measuring the quality of code
translation systems. BLEU remains useful as a fast and simple similarity-based
metric, but it should be used together with execution-based methods rather than
as the only evaluation criterion.